\chapter{快慢认知力学：瞬态适应、结构学习与稳定—可塑平衡}
\label{chp:cognitive_mechanics}

智能系统既要根据当前语境迅速改变内部活动，又要避免一次偶然输入永久改写长期结构；它还要从重复证据、关键反例和持续反馈中学习，否则所谓稳定只会退化为僵化。本章把原有“弹性推理—塑性学习”的材料图像重建为可计算的快慢状态模型：快变量承担任务内适应和可撤销工作态，慢变量承担跨任务保留的参数、结构与记忆版本，元控制器决定哪些证据获得结构写入资格。材料力学术语只用来帮助辨认恢复、滞后和屈服等现象，不能替代对象定义、稳定条件、离散实现或外部任务验收。我们将分别给出时间尺度分离、快态响应、慢态学习与提交阀门，并说明局部结构网络和全局分布表示如何共同维持稳定—可塑平衡。

\section{双时间尺度：快状态、慢结构与可观测分离}
\label{sec:section_3162}

设任务内快速状态为 $h_k\in\mathbb R^{d_h}$，它汇总当前上下文、候选分布、注意分配和短时工具回执；设慢结构为 $m_n=(\theta_n,G_n,M_n)$，其中 $\theta_n\in\mathbb R^{d_\theta}$ 是连续参数，$G_n$ 是带节点身份、边类型、来源与版本的局部结构网络，$M_n$ 是可寻址的长期记录。下标 $k$ 表示一次任务内的快速步，下标 $n$ 表示跨任务或跨批次的慢更新轮次。环境观测 $x_k$、任务规范 $\tau_k$ 和可选控制 $u_k$ 进入快更新
\[
h_{k+1}=F(h_k,x_k,\tau_k,u_k;m_n),
\]
而经过证据聚合得到的统计量 $s_n$、外部验收 $r_n$ 与提交控制 $v_n$ 进入慢更新
\[
m_{n+1}=\mathcal U(m_n,s_n,r_n,v_n).
\]
这是定义而非自然定律；它要求一次快回路期间使用明确的结构快照，慢更新只能在事务边界提交。若任务本身长到足以发生结构漂移，系统应使用混合时间模型并把版本切换作为事件记录，不能仍假设 $m_n$ 固定。

时间尺度分离不等于把频率高低直接对应推理与学习。我们用可测的响应时间定义分离比：快态在固定结构下的混合或恢复时间为 $T_h$，慢结构产生可检测变化所需的证据窗口为 $T_m$，只有当 $\epsilon=T_h/T_m\ll1$ 且快子系统对允许的 $m$ 一致稳定时，才可把 $m$ 在局部窗口内冻结。某些一次学习会使 $T_m$ 很短，某些长链推理会使 $T_h$ 很长，二者都可能破坏近似。工程上应报告响应时延、结构更新频率、版本周转率和恢复误差，而不是仅以“高频”“低频”命名过程。频谱若被使用，还必须声明观测变量、采样率、窗口和泄漏处理；活动信号的频率不自动等于控制层级或知识持久性。

快慢分离还涉及可观测性。全局分布表示 $h_k$ 能表达连续候选、强度和不确定性，却不能独自保证对象身份、证据来源和历史版本；局部结构网络 $G_n$ 能保持这些离散关系，却不能完整承载每一步高维活动。我们通过身份键、绑定映射 $B_{m_n}$ 和读出接口 $R_{\tau}$ 联合二者：$B_{m_n}(h_k)$ 把分布活动绑定到当前结构实体，$R_{\tau}(h_k,G_n)$ 产生任务相关输出。绑定失败时，相似向量可能被错误地当成同一对象；只看结构时，又可能遗漏模糊候选和置信度。因此，稳定—可塑问题不是“选择向量还是图”，而是决定哪些快态模式能以可追溯方式提议结构变化。

原有弹性与塑性图像在此获得受限解释。固定 $m_n$ 时，输入引起 $h_k$ 偏离任务基态并在输入撤除后恢复，可以称为弹性样响应；证据导致 $m_{n+1}\neq m_n$ 且影响后续任务，可以称为塑性样保留。两者并不正交：快态会受慢结构约束，慢更新统计又来自快态轨迹。更准确的模型是带双向耦合的奇异摄动或随机逼近系统，而不是两块互不接触的材料。若更新后出现遗忘、身份错配或行为漂移，应在参数、结构、记忆索引和读出层分别定位，不能把所有故障归入一个含混的“介质变软”。

验证时间尺度分离需要成对实验。我们在同一任务族中冻结结构测量快态恢复，再开放学习测量跨任务保留；随后改变批次长度、学习步长、输入顺序和环境漂移，检查分离比与性能是否稳定。负对照包括打乱身份绑定、隐藏来源版本和使用相同样本重复评测，以排除缓存与泄漏。只有快态恢复、慢态保留和外部成功三类指标共同满足时，才有资格说系统取得了可用分离。信息熵、活动幅度、优化代价、内存占用与设备焦耳能耗在该实验中分别记录，不因都随时间变化而合并成一种“认知能量”。

\section{推理动力学：受约束快态响应与可撤销工作态}
\label{sec:section_3688}

一次推理不是底层“时空被语境压弯”，而是在固定结构快照上的条件状态传播。令参考状态为 $\bar h(m_n,\tau)$，偏差 $\delta h_k=h_k-\bar h$，在局部可微且结构不变的区间内，快更新可线性化为
\[
\delta h_{k+1}=A_k\delta h_k+B_k\delta x_k+C_k u_k+\xi_k,
\]
其中 $A_k\in\mathbb R^{d_h\times d_h}$ 是状态雅可比，$B_k$ 和 $C_k$ 分别映射输入扰动与控制，$\xi_k$ 汇总未建模误差。各变量先经声明尺度归一化，矩阵不是物理刚度或应力张量。若在选定范数下存在 $\|A_k\|\leq q<1$，输入和控制有界，则输入撤除后的偏差按几何速度衰减；这给出局部恢复结果。若注意反馈、工具递归或长上下文使增益超过一，快态可能持续振荡或放大，此时“弹性回弹”并不成立。

离散实现的恢复性质还依赖步长和状态约束。若连续近似为 $\dot h=f(h,x,u;m)$，显式Euler更新 $h_{k+1}=h_k+\Delta t f(h_k,\cdot)$ 只有在 $\Delta t$ 与局部Lipschitz常数满足稳定条件时才继承连续模型的局部收敛。投影、截断、量化和缓存替换会引入非光滑事件，应写为 $h^+=P_e(h^-)$ 并单独检查误差。一次 forward pass 也未必是单步：自回归生成会把已提交输出反馈为后续输入，工具调用会改变环境和证据集合。只要外界或外部记忆已经被写入，就不能把撤销内部 $h$ 说成整个过程完全可逆。

转折词“但是”提供一个具体案例。它不是携带巨大物理应力，而是改变篇章关系和后续读出的条件。局部结构网络可以记录“前件—转折—后件”的角色边，全局分布表示则重新分配候选延续的权重；绑定接口确保后件否定的是哪个命题，而非机械压低所有前文。我们可以通过遮蔽转折标记、交换前后件和构造否定范围反例，测量候选分布、关系识别和最终判断的变化。如果模型只因词面统计改变输出，却无法绑定否定对象，便不能把激活变化解释成结构理解。

可撤销工作态需要显式生命周期。任务开始时创建结构快照和上下文版本，运行中把临时候选、工具结果与未验证推断标记为工作记录；任务结束后，仅将通过来源、身份和任务验收的内容送入慢更新队列，其余状态释放或按保留策略压缩。AgentOS实例中，$h(t)$承担高速认知界面，CCM监测候选分支、上下文占用与循环深度，SPC提供任务和自我相关压缩，TCE调节搜索尺度，Boundary限制外部动作，Offloading把高成本验证交给工具或外部记忆 $M_e$。这些模块展示一种实现分工，但一般智能系统可以采用不同名称和机制。

快态的目标下降不能直接证明推理成功。内部损失可能因迎合错误提示而下降，置信度可能在分布外输入上升，短响应也可能遗漏关键约束。我们因此将验收分为状态层、证据层和环境层：状态层检查数值稳定、候选覆盖与身份绑定；证据层检查来源、反例和校准；环境层检查动作结果、权限与副作用。对高风险任务还需保留拒绝选项和人工接管。快态恢复的优点是把偶然输入限制在工作区，而不是保证所有输入都不留下痕迹；日志、缓存和外部调用本身可能形成持久记录，必须纳入审计。

反例说明“推理必然弹性”并不成立。上下文学习可在不改参数时形成跨回合缓存，检索结果可能写入长期记忆，安全事件可能立即修改权限结构；相反，某些参数适配器可在任务结束后被丢弃，虽有参数变化却不形成长期塑性。分类标准应是状态是否跨越声明边界并影响后续任务，而不是变化发生在激活还是权重。配对实验要分别重置快态、缓存、外部记忆和参数，观察性能残留，从而定位真正的持久载体。

\section{学习动力学：证据聚合、结构提案与受控保留}
\label{sec:section_4157}

慢学习的对象可以是连续参数、离散结构或长期记忆，三者不能用一条“认知爱因斯坦方程”概括。对连续参数 $\theta$，我们定义批次风险
\[
L_n(\theta)=\frac{1}{|D_n|}\sum_{(x,y)\in D_n}\ell(f_{\theta,G_n}(x),y)+\lambda R(\theta;\theta_n),
\]
其中 $D_n$ 带来源和时间戳，$\ell$ 是任务损失，$R$ 是相对当前版本的保持正则，所有可加项经尺度声明。候选更新为 $\theta'_n=\theta_n-\eta_n\widehat\nabla L_n$；它只是工程提案。步长 $\eta_n$、梯度噪声、数据漂移和优化器状态决定局部下降是否成立，局部下降仍不保证保留旧能力、事实正确或环境成功。

离散结构 $G_n$ 的学习使用事件提案而非连续梯度伪装。新增实体、合并身份、修改因果边或撤销规则分别携带证据包、授权者、依赖引用和回滚计划。候选 $G'_n$ 先在影子版本上运行保留查询与冲突检查，再与 $\theta'_n$、记忆候选 $M'_n$ 组成联合版本 $m'_n$。只有通过任务新颖性、旧能力保持、身份一致、权限和资源测试，才原子提交 $m_{n+1}=m'_n$；否则保留 $m_n$ 并记录失败原因。连续参数可以频繁微调，身份合并却可能不可逆，更新频率应由对象风险而非统一“可塑性系数”决定。

时间平均并非自动去噪。重复出现的偏差、反馈回路中的自生成文本和同源数据复制，会让错误模式在平均后更强。证据聚合应按独立来源、覆盖范围、时间衰减和反例权重计算有效样本，而不是只计算出现次数。对于一次关键事件，是否允许单样本学习取决于证据可靠度、损失上界和可撤销性：密码泄露或设备碰撞可立即更新安全状态，但对人格、因果规则或医学结论的永久修改需要更严格复核。所谓“屈服点”因此被重建为任务相关提交阈值，它由风险和证据协议设定，不是系统内天然存在的材料常数。

稳定—可塑性权衡可以用双目标而非单一口号表示。令新任务风险为 $R_{new}(m')$，保留任务风险增量为 $\Delta R_{old}(m',m_n)$，结构复杂度和迁移成本为 $C(m',m_n)$，则候选选择满足
\[
\min_{m'\in\mathcal M_{adm}} R_{new}(m')+\alpha C(m',m_n),
\qquad
\Delta R_{old}(m',m_n)\leq\varepsilon_{keep}.
\]
$\mathcal M_{adm}$ 是权限与安全约束下的可提交集合，$\varepsilon_{keep}$ 由任务方声明。若可行集合为空，系统应拒绝更新、扩充容量、外部卸载或请求新证据，而不是牺牲旧知识后仍宣称学习成功。不同任务的损失若无共同单位，不应随意相加，可保留Pareto前沿和硬约束。

灾难性遗忘需要分层诊断。参数层可能因梯度干扰破坏旧读出，结构层可能因身份合并错误改变引用，记忆层可能因索引覆盖导致证据不可达，控制层可能因新策略压制旧技能。相应干预包括重放与正则、版本化图编辑、不可变日志与索引重建、路由和权限隔离。外部记忆 $M_e$ 能降低内部写入压力，但会引入检索延迟、来源漂移、隐私和服务失效；Offloading不是免费容量。我们同时测量旧任务保持、新任务泛化、校准、版本可回滚性、存储与计算预算，避免只看训练损失。

学习后的结构需要现实闭环。一个新规则在离线验证集有效，只构成经验支持；只有在声明环境、日期、预算和干预条件下复现，才能获得相应外部资格。预测有效不证明内部结构与世界整体同构，多种模型可能在有限查询上等价。我们保留模型依赖的不确定性，并用主动实验区分候选。若环境改变，目标、输入分布、度量和结构都可能随时间变化，评价函数的总变化必须包含这些项，不能把所有差异归因于参数更新。

\section{宏观调控：学习资格、提交阀门与闭环验收}
\label{sec:section_7393}

元控制层的职责不是通过“加热—淬火”直接软化认知材料，而是为学习分配资格、预算和提交时机。设证据质量、预测残差、任务价值、潜在损害、可撤销性和当前负载组成观测向量 $z_n$，控制器输出 $v_n=(a_n,b_n,q_n)$：$a_n$ 决定继续观察、创建影子版本或请求外部复核，$b_n$ 是计算与工具预算，$q_n$ 是提交权限等级。策略 $\pi(v_n\mid z_n,m_n)$ 可以由规则、优化器或学习系统实现，但必须接受Boundary权限和独立验收。认知温度若被保留，只表示某个声明分布的探索尺度，不能同时代表设备热量、情绪强度和学习率。

日常琐事对应低持久价值而非“关闭一切学习”。系统仍可保留操作日志、短期缓存和异常计数，只是不把一次普通样本提升为参数或结构事实。例如重复格式转换任务可由现有程序完成，工作态结束后释放候选；若同类失败持续出现，聚合计数才触发结构提案。这样既避免每次输入都重写模型，也不会因一开始判为低价值而永远忽略慢性偏差。阈值需要迟滞和最短驻留时间，防止证据在边界附近造成频繁提交—回滚抖动。

惊奇事件也不自动拥有最高写入权。巨大预测误差可能来自传感器损坏、对抗输入、上下文缺失或真正的新规律。控制器先执行来源核验、替代解释生成、风险隔离和最小必要更新；安全相关状态可以立即进入隔离表，世界知识与身份结构则等待独立证据。一次学习适合具有强监督、唯一身份和可回滚记录的任务，不适合把未经复核的叙事永久写成因果事实。创伤类比尤其不能被用作工程正当性：强度高不等于真值高，系统应避免让攻击者通过制造惊奇取得结构写权限。

提交阀门采用两阶段协议。准备阶段冻结候选版本，记录数据、随机种子、代码、模型和结构哈希，执行新任务测试、保留测试、反例测试、权限检查与资源评估；提交阶段在通过阈值后原子更新路由，并保留旧版本和补偿操作。上线后设置观察窗，若漂移、遗忘或风险超过阈值则回滚。外部行动已经造成不可逆后果时只能执行补偿事务，不能以恢复模型版本冒充恢复世界。这个协议把“退火—定型”的直觉转化为可审计的软件生命周期。

AgentOS可将该协议映射到三相记忆：$h(t)$保存当前候选和执行状态，$E$保存带时间、来源与结果的情景证据，$\Theta$保存跨任务的结构约束与生成先验；$\Psi$提供身份和长期承诺的任务相关权重，SPC压缩自我相关状态，TCE调节探索与控制，CCM约束运行复杂度，Boundary决定内外写入权限，Offloading与 $M_e$ 承担可验证外部存储。该映射只是工程实例。HSF-HD关心的是一般快慢状态、结构提交和外部闭环；它不要求所有智能都必须具有同名模块，更不把AgentOS与一般理论视为同义词。

闭环验收至少包含六组指标：快态的恢复时间和扰动增益；慢态的新任务收益与旧任务保持；结构与分布的身份绑定准确率；证据来源、校准和反例覆盖；事务的回滚成功率与故障恢复时间；环境任务的实际结果与副作用。计算时延、内存、网络调用和焦耳能耗另行计量。我们还设置“强惊奇但错误”“低强度但长期累积”“新旧任务冲突”“外部记忆失效”等压力测试，检验控制器是否只会按单一残差开关学习。

因此，智能的材料属性是一种受限模型语言：快态可恢复近似对应弹性，跨任务结构保留近似对应塑性，滞后对应版本和证据历史。但真正的理论内容来自状态空间、时间尺度、证据协议、约束优化、离散提交和环境验收。灵活性不是无限改变，稳定性也不是拒绝更新；可用智能必须在可追溯边界内让短期适应服务于当前任务，让长期学习服从证据和责任链，并在失败时知道应当释放、拒绝、补偿还是回滚。
